\workbookchapter{领域应用}

\begin{exerciselist}
\exerciseitem 将“提升客服效率”改写为可计数任务，定义输入、输出、结束条件与失败。

\begin{workbookanswers}
可定义“在授权知识范围内处理一条退换货咨询”：输入订单身份、用户问题及规则版本；输出附证据答复或需人工原因；结束为用户所需信息已核验给出或明确转交。记录每请求处理时间、首次解决、错误承诺和人工干预，超时也是失败状态。满意度可附带，但不是唯一结束判据。
\end{workbookanswers}

\exerciseitem\optional 为合同条款抽取设计证据记录，说明为何仅保存文件名不足以复核。

\begin{workbookanswers}
证据记录包含合同ID与哈希、修订、生效日、页码及条款位置、原文跨度、结构化字段、单位与例外、提取器版本和核验人。文件名可以重复或内容被替换，不能固定证据对象。把“自动续期”提取为布尔值时还须保存期限、通知条件及例外，否则结构虽正确却丢失法律语义。
\end{workbookanswers}

\exerciseitem\optional 对库存查询列出身份、单位、状态与时间字段，并构造遗漏任一字段导致错误的例子。

\begin{workbookanswers}
至少包含租户/查询者、SKU及仓库、数量单位、在库/可用/冻结等状态、业务时间和快照版本。缺身份可查到其他租户；缺单位把箱数当件数；缺状态把冻结库存当可用；缺时间把昨日余额当当前承诺。涉及实际预留还需原子条件写，查询正确不保证稍后库存仍足。
\end{workbookanswers}

\exerciseitem 根据式\tbobject{eq:domain-timesave}计算辅助操作5分钟、复核4分钟、30\%任务返工8分钟时，相对12分钟原流程的净收益。

\begin{workbookanswers}
新期望时间为$5+4+0.3\times8=11.4$分钟，净节省$12-11.4=0.6$分钟/任务，即5\%。100任务只节省60分钟，维护、评估和系统成本未计；若这些摊销超过0.6分钟等价成本则净收益可变负。不能仅比较原12与辅助5得7分钟节省。
\end{workbookanswers}

\exerciseitem 将自动与人工两种决策都加入执行成本及残余错误，重新推导自动化阈值。

\begin{workbookanswers}
令自动成本$a$、正确概率p、错损L；人工成本$h$、残余错率e，暂设错损相同且正确损失为0。自动风险$a+(1-p)L$，人工风险$h+eL$；自动优于人工要求 $p>1-e+\frac{a-h}{L}$，L须正，阈值超范围时作总选某行动解释。若人工错误也依案例变化，使用$e(x)$并逐例比较；拒答等第三动作须一并纳入。
\end{workbookanswers}

\exerciseitem\optional 为制造业告警分析区分观察、统计关联、故障假设与已确认原因。

\begin{workbookanswers}
观察是“14:02温度传感器读数超阈值”；关联是“历史同类告警常伴随冷却流量下降”；假设是“可能堵塞导致冷却不足”；确认原因需现场检查、控制实验或维护证据证明具体堵塞并排除替代解释。模型可组织证据和提出假设，不能将相关性自动升级为已确认故障，更不能据此绕过工业控制许可。
\end{workbookanswers}

\exerciseitem\optional 为什么教育对话满意度不能直接证明学习效果？设计更贴近目标的测量。

\begin{workbookanswers}
满意度受流畅、迎合和即时情绪影响，学习效果要观察知识或迁移技能的变化。可按学生随机分组或分层，使用等价前测后测、延迟保持测试和未见过的迁移题，控制教师与学习时长，按学生为独立单位。需要报告原始水平、缺失和帮助次数，不能把多轮点赞当独立学习增益。
\end{workbookanswers}

\exerciseitem\optional 给出一项代码任务，说明编译、测试、评审与真实部署分别证明什么。

\begin{workbookanswers}
以修复金额舍入为例：编译仅证明语法和类型满足编译器；单元/集成测试检查有限输入及接口；评审检查设计、边界和可维护性但仍有漏检；部署验证在真实依赖和流量中是否满足指标。每层应保存对应证据，测试通过不代表已发布，发布成功也不自动证明所有金额正确。
\end{workbookanswers}

\exerciseitem 构造上线前后成功率改善但实际处理能力未改善的任务分布变化。

\begin{workbookanswers}
简单任务正确率0.95、困难0.50固定。上线前各占一半，总体0.725；上线后简单占80\%，总体0.86。总分上升0.135而各切片能力未变。应以固定任务分布加权比较，另报告真实流量构成，避免把更容易的请求归功于模型升级。
\end{workbookanswers}

\exerciseitem\optional 设计以设备为分组单位的评估划分，解释随机划分日志行的泄漏风险。

\begin{workbookanswers}
将同设备全部日志放同一训练/验证/测试组，测试设备不能在训练中出现；若预测未来，还按时间留出。随机按行会让相邻传感器窗口、设备ID和同一次故障模式泄入两侧，结果可能依赖记忆而非跨设备泛化。应另测新设备和已知设备未来两个目标，不混称同一个泛化问题。
\end{workbookanswers}

\exerciseitem 在自动覆盖算例中，将所有自动错误交给人工返工，重新定义干预率与成本分母。

\begin{workbookanswers}
原1000任务中300经人工，700自动含70错误。若70错误都被识别并转人工，则至少370个不同任务受干预，干预率37\%；700不是最终无人工覆盖，真正完全自动且正确仅630，即63\%。新增70的返工成功率与耗时题设未给，不能断言最终成功970；若其成功率r，则总成功$900+70r$。成本应以1000全部提交任务平均，含自动尝试后返工双重成本。
\end{workbookanswers}

\exerciseitem 为一个领域应用制定可复现失败记录，明确哪些敏感信息不需要保存原文。

\begin{workbookanswers}
保留任务类别、版本化输入引用、目标与实际行为、证据版本、工具动作/回执、终止原因及最小复现条件。敏感正文可放受控存储并用短期授权引用，日志保存哈希、字段类别或脱敏片段；凭据与无关个人信息不应记录。删除/保留期限、访问审计和回归样本用途需清楚，失败可复现不等于复制全部生产数据。
\end{workbookanswers}

\exerciseitem 制造业知识助手根据维修文档给出高置信故障判断，并建议调整设备控制参数。设计从证据检索到允许执行动作的流程，区分模型置信度、操作权限和独立安全约束，并说明审计记录应包含哪些信息。

\begin{workbookanswers}
先核对文档版本、设备型号与运行条件，区分观测事实、相关线索、诊断假设和已确认原因。模型置信度可作为待评估信号，不能授予操作权限或证明物理安全。业务动作应由受控接口校验身份、对象、参数范围和前置状态，按风险执行审批，并由独立的确定性安全约束与联锁限制危险控制。记录证据来源、建议、批准主体、执行前状态、实际参数、执行结果及失败处置；对重复请求与状态变化规定幂等和一致性策略。执行后的独立状态核验用于确认效果，模型生成的成功描述不能代替设备证据。
\end{workbookanswers}

\end{exerciselist}
